\section{Results}
\label{sec:results}

All numerical claims below are taken from verified artifacts \texttt{results/manuscript\_ready\_numbers.json}, \texttt{results/v3\_gated/t1\_architecture\_selection.json}, and \texttt{results/aggregate\_v3.json} (baselines/T2/RCA), without modifying frozen instances.

\subsection{T1 anomaly and T2 failure prediction}
Table~\ref{tab:t1t2} summarizes multi-seed mean ranking metrics. On T1, the final ECN-v3 anchored detector achieves mean AUPRC $0.1152$ (bootstrap 95\% CI $[0.0679,0.1680]$; architecture-selection artifact \texttt{results/manuscript\_ready\_numbers.json}), exceeding random forest ($0.0758$) and the v2 legacy configuration ($0.0577$). Mean T1 ROC-AUC is $0.7534$. Fig.~\ref{fig:baseline} visualizes the same AUPRC means with 95\% confidence intervals for T1 and T2, highlighting the T1 lift of the final head and the T2 telem-logistic recommendation.

\begin{table*}[!t]
\centering
\caption{Multi-seed mean detection metrics on T1 (anomaly) and T2 (failure horizon). ECN-v3 uses leakage-safe enriched features with \textbf{anchored} fusion on T1; the recommended T2 head is telemetry logistic regression. AUPRC entries report means over six frozen instances with 95\% confidence intervals.}
\label{tab:t1t2}
\renewcommand{\arraystretch}{1.18}
\setlength{\tabcolsep}{3.5pt}
{\footnotesize
\begin{tabular*}{\textwidth}{@{\extracolsep{\fill}}lcccc}
\toprule
Method & T1 AUPRC & T1 ROC-AUC & T2 AUPRC & T2 ROC-AUC \\
\midrule
\textbf{ECN-v3 (final)$^{\dagger}$} & $\mathbf{0.1152}$ $[0.0609,\,0.1695]$ & $\mathbf{0.7534}$ & $0.0380$ $[-0.0027,\,0.0788]$ & $0.6891$ \\
Stacking ablation & $0.0997$ $[0.0738,\,0.1257]$ & $0.7803$ & --- & --- \\
v2 legacy+anchored & $0.0577$ $[0.0214,\,0.0940]$ & $0.6831$ & --- & --- \\
Random forest & $0.0758$ $[0.0505,\,0.1011]$ & $0.7661$ & $0.0176$ $[0.0099,\,0.0253]$ & $0.7081$ \\
Logistic & $0.0631$ $[0.0329,\,0.0934]$ & $0.7260$ & $\mathbf{0.0380}$ $[-0.0027,\,0.0788]$ & $0.6891$ \\
LightGBM & $0.0570$ $[0.0335,\,0.0805]$ & $0.6798$ & $0.0213$ $[-0.0041,\,0.0467]$ & $0.5483$ \\
Isolation Forest & $0.0575$ $[0.0331,\,0.0819]$ & $0.6716$ & $0.0216$ $[-0.0040,\,0.0472]$ & $0.6459$ \\
EWMA & $0.0403$ $[0.0198,\,0.0607]$ & $0.5484$ & $0.0066$ $[0.0045,\,0.0086]$ & $0.4482$ \\
GNN proxy & $0.0335$ $[0.0167,\,0.0504]$ & $0.6846$ & $0.0106$ $[0.0069,\,0.0144]$ & $0.5965$ \\
MLP sequence & $0.0084$ $[0.0053,\,0.0115]$ & $0.3617$ & $0.0054$ $[0.0041,\,0.0067]$ & $0.3795$ \\
Majority & $0.0098$ $[0.0083,\,0.0113]$ & $0.5000$ & $0.0056$ $[0.0050,\,0.0061]$ & $0.5000$ \\
\bottomrule
\multicolumn{5}{l}{\scriptsize $^{\dagger}$T1: \texttt{ECNFusionModel} on v3 features; T2 column reports the recommended telem-LR head (identical to Logistic on T2).} \\
\multicolumn{5}{l}{\scriptsize Stacking and v2 rows are historical/ablation references, not the final proposed detector.}
\end{tabular*}}
\end{table*}


\begin{figure}[!t]
\centering
\includegraphics[width=\linewidth]{baseline_auprc_mean_ci.pdf}
\caption{Multi-seed mean AUPRC with 95\% confidence intervals. Left: T1 anomaly detection, where ECN-v3 denotes the final anchored head on enriched features. Right: T2 failure-horizon ranking; the ECN T2 column is the recommended telemetry logistic head.}
\label{fig:baseline}
\end{figure}

On T2, the recommended head is telemetry logistic regression with mean AUPRC $0.038043099063649714$, higher than RF ($0.017610244849363122$). We do not claim that multi-agent fusion improves T2 under ultra-rare positives.

Paired Wilcoxon tests and Cliff's $\delta$ (Table~\ref{tab:sig}) show that the T1 mean advantage versus random forest is not significant at $\alpha{=}0.05$ ($p{=}0.688$, $\delta{=}{+}0.389$), consistent with $n{=}6$ underpowering. Versus logistic and several weaker baselines, paired $p$-values reach $0.031$. Stacking versus anchored is not significant ($p{=}0.375$).

\begin{table*}[!t]
\centering
\caption{Paired Wilcoxon signed-rank tests and Cliff's $\delta$ for T1 AUPRC of the final ECN-v3 anchored model versus baselines and the stacking ablation ($n{=}6$). Positive $\delta$ favors ECN-v3.}
\label{tab:sig}
\renewcommand{\arraystretch}{1.12}
\setlength{\tabcolsep}{4pt}
{\footnotesize\begin{tabular*}{\textwidth}{@{\extracolsep{\fill}}lcccc}
\toprule
Comparator & $p$ & Cliff $\delta$ & ECN-v3 mean & Comparator mean \\
\midrule
Random forest & $0.688$ & $+0.389$ & $0.1152$ & $0.0758$ \\
Logistic & $0.031$ & $+0.556$ & $0.1152$ & $0.0631$ \\
LightGBM & $0.062$ & $+0.611$ & $0.1152$ & $0.0570$ \\
Isolation Forest & $0.031$ & $+0.667$ & $0.1152$ & $0.0575$ \\
EWMA & $0.062$ & $+0.778$ & $0.1152$ & $0.0403$ \\
GNN proxy & $0.062$ & $+0.778$ & $0.1152$ & $0.0335$ \\
MLP sequence & $0.031$ & $+1.000$ & $0.1152$ & $0.0084$ \\
Majority & $0.031$ & $+1.000$ & $0.1152$ & $0.0098$ \\
Stacking ablation & $0.375$ & $+0.111$ & $0.1152$ & $0.0997$ \\
\bottomrule
\end{tabular*}}
\end{table*}


Representative ROC and precision--recall curves on the primary frozen instance appear in Fig.~\ref{fig:t1curves}. Curve geometry remains consistent with rare-positive ranking: ROC discrimination is moderate, while precision--recall is challenging for all methods. Caption and legend identify the final \textbf{anchored} ECN-v3 head; any historical stacking traces in stored JSON are ablations only and are not used for mean claims.

\begin{figure}[!t]
\centering
\includegraphics[width=\linewidth]{T1_v1.1.0-INST_roc_pr.pdf}
\caption{T1 ROC (left) and precision--recall (right) curves on frozen instance v1.1.0-INST for representative detectors. Legend values report ROC-AUC and average precision on this instance.}
\label{fig:t1curves}
\end{figure}

\subsection{Architecture selection and fusion ablation}
Fig.~\ref{fig:archsel} and Table~\ref{tab:ablation} isolate the sources of the T1 improvement. Leakage-safe feature enrichment accounts for the dominant gain; replacing anchored fusion with stacking on the same enriched features \emph{reduces} mean AUPRC from $0.1152$ to $0.0997$. Accordingly, stacking is reported as an ablation/negative result rather than as the final T1 architecture.

\begin{figure}[!t]
\centering
\includegraphics[width=0.92\columnwidth]{architecture_selection_t1.pdf}
\caption{T1 architecture selection: v2 legacy+anchored, v3+stacking (ablation), v3+anchored (final), and RF telem\_only. Error bars denote 95\% confidence intervals over six seeds.}
\label{fig:archsel}
\end{figure}

\begin{table}[!t]
\centering
\caption{Ablation mean AUPRC on T1/T2 (six seeds, 95\% CI).}
\label{tab:ablation}
\renewcommand{\arraystretch}{1.1}
\setlength{\tabcolsep}{3.5pt}
{\footnotesize\begin{tabular}{llcc}
\toprule
Task & Ablation & Mean & 95\% CI \\
\midrule
T1 & full & $0.111$ & $[0.042,\,0.180]$ \\
 & no\_twin & $0.115$ & $[0.043,\,0.187]$ \\
 & no\_nbr & $0.116$ & $[0.051,\,0.181]$ \\
 & telem\_only & $0.069$ & $[0.014,\,0.123]$ \\
 & twin\_only & $0.028$ & $[0.013,\,0.044]$ \\
T2 & full & $0.015$ & $[0.004,\,0.026]$ \\
 & no\_twin & $0.014$ & $[0.003,\,0.025]$ \\
 & no\_nbr & $0.014$ & $[0.003,\,0.025]$ \\
 & telem\_only & $0.036$ & $[-0.007,\,0.078]$ \\
 & twin\_only & $0.016$ & $[0.008,\,0.023]$ \\
\bottomrule
\end{tabular}}
\end{table}


Table~\ref{tab:modules} and Fig.~\ref{fig:modules} summarize verified contribution deltas. Feature enrichment versus v2 contributes ${+}0.0575$ mean AUPRC. Twin contribution under the final head is near zero (${+}0.0003$). Stacking relative to anchored contributes ${-}0.0155$.

\begin{table}[!t]
\centering
\caption{Module AUPRC contribution ($\mathrm{full}-\mathrm{ablated}$) from full-suite ablations.}
\label{tab:modules}
\renewcommand{\arraystretch}{1.1}
{\footnotesize\begin{tabular}{lccc}
\toprule
Task & Telemetry & Neighbor & Digital Twin \\
\midrule
T1 & $0.082$ & $-0.005$ & $-0.004$ \\
T2 & $-0.000$ & $0.001$ & $0.001$ \\
\bottomrule
\end{tabular}}
\end{table}


\begin{figure}[!t]
\centering
\includegraphics[width=0.78\columnwidth]{module_contribution.pdf}
\caption{Verified T1 AUPRC deltas: feature enrichment versus v2, twin contribution under the final head, and stacking versus anchored fusion.}
\label{fig:modules}
\end{figure}

\subsection{Calibration and operating points}
Calibration quality is summarized in Table~\ref{tab:cal}. Final T1 raw Brier is $0.1509$, worse than RF ($0.0227$), motivating optional Beta calibration for probability displays without altering AUPRC ranking. Fig.~\ref{fig:t1cal} shows a reliability diagram on v1.1.0-INST for stored ECN scores.

\begin{table}[!t]
\centering
\caption{Calibration quality (Brier score; lower is better). Final ECN-v3 T1 Brier is reported from the architecture-selection study; baseline Brier values are from the six-seed aggregate. Optional Beta calibration is recommended for probability display and does not change AUPRC ranking.}
\label{tab:cal}
\renewcommand{\arraystretch}{1.15}
{\footnotesize\begin{tabular}{llcc}
\toprule
Task & Method & Mean & Notes \\
\midrule
T1 & ECN-v3 final (raw) & $0.1509$ & Optional Beta for display \\
T1 & Random forest & $0.0227$ & Strong tabular calibration \\
T1 & Logistic & $0.2328$ & \\
T2 & Telem logistic (T2 head) & $0.2577$ & Recommended T2 head \\
T2 & Random forest & $0.0124$ & \\
\bottomrule
\end{tabular}}
\end{table}


\begin{figure}[!t]
\centering
\includegraphics[width=0.78\columnwidth]{T1_v1.1.0-INST_calibration.pdf}
\caption{T1 reliability diagram on instance v1.1.0-INST. The dashed line denotes perfect calibration. Optional Beta calibration is recommended for operator-facing probabilities.}
\label{fig:t1cal}
\end{figure}

\subsection{RCA, healing decision support, and cost}
Table~\ref{tab:rcaheal} reports RCA macro-F1, healing decision-support scores, and cost. Proposed T3 RCA macro-F1 mean is $0.25396825396825395$ using RF with TreeSHAP explanations. Healing with RCA category features reaches macro-F1 $1.0$ on this synthetic setup---an upper bound that is \emph{not} interpreted as operational perfection; the honest \texttt{no\_rca\_cat} setting yields $0.19027777777777777$. Fig.~\ref{fig:t3cm} illustrates a representative T3 confusion matrix on the primary instance.

\begin{table}[!t]
\centering
\caption{RCA macro-F1, healing scores, and evaluation wall-clock cost (six seeds).}
\label{tab:rcaheal}
\renewcommand{\arraystretch}{1.1}
{\footnotesize\begin{tabular*}{\columnwidth}{@{\extracolsep{\fill}}llc}
\toprule
Metric & Setting & Mean [95\% CI] \\
\midrule
T3 RCA macro-F1 & proposed & $0.254$ $[0.090,\,0.418]$ \\
Healing macro-F1 & with RCA category & $1.000$ $[1.000,\,1.000]$ \\
Healing macro-F1 & without RCA category & $0.190$ $[0.099,\,0.282]$ \\
Eval.\ wall time (s/seed) & full suite & $119.7$ $[111.8,\,127.6]$ \\
\bottomrule
\end{tabular*}}
\end{table}


\begin{figure}[!t]
\centering
\includegraphics[width=0.82\columnwidth]{T3_v1.1.0-INST_cm.pdf}
\caption{Representative T3 RCA confusion matrix for the RF+TreeSHAP RCA head on instance v1.1.0-INST.}
\label{fig:t3cm}
\end{figure}

\subsection{Why simplicity remains preferable}
At $\sim$1\% positives, enriched telemetry recovers most ranking signal. Anchored fusion is simpler than stacking, improves mean T1 AUPRC on the same features, and remains competitive with RF in mean while acknowledging non-significant paired tests versus RF. With $n{=}6$ seeds, intervals are wide; claims emphasize reproducible mean improvements and architecture honesty rather than universal paired dominance. Deep tabular (TabNet) and true GraphSAGE message passing did not displace RF or ECN-v3 on mean T1 AUPRC under this protocol (Section~\ref{sec:deep_baselines}).

\subsection{Deep tabular and true GNN baselines}
\label{sec:deep_baselines}
Table~\ref{tab:deep_baselines_v5} reports measured six-seed means for TabNet and true GraphSAGE, alongside the historical GNN proxy and RF. Artifacts: \texttt{results/aggregate\_v3.json} keys \texttt{tabnet\_\_full}, \texttt{graphsage\_\_full}, \texttt{gnn\_graphsage\_proxy\_\_full}; manuscript pointer \texttt{extensions\_v5\_deep\_baselines} in \texttt{results/manuscript\_ready\_numbers.json}.

\begin{table}[!t]
\centering
\caption{Deep baselines on frozen six-seed T1/T2 (mean AUPRC with 95\% CI). TabNet uses telem\_only features; true GraphSAGE performs mean-aggregation message passing on the Digital Twin device graph. The historical GNN proxy remains LightGBM (not message-passing). Final ECN-v3 T1 selection numbers remain those in \texttt{manuscript\_ready\_numbers.json}.}
\label{tab:deep_baselines_v5}
\renewcommand{\arraystretch}{1.15}
{\footnotesize\begin{tabular}{lcc}
\toprule
Method & T1 AUPRC & T2 AUPRC \\
\midrule
ECN-v3 final (manuscript) & $0.1152$ $[0.0609,0.1695]$ & --- \\
TabNet & $0.0492$ $[0.0213,0.0772]$ & $0.0110$ $[0.0038,0.0183]$ \\
GraphSAGE (true GNN) & $0.0152$ $[0.0059,0.0246]$ & $0.0062$ $[0.0057,0.0068]$ \\
GNN proxy (LightGBM) & $0.0335$ $[0.0167,0.0504]$ & $0.0106$ $[0.0069,0.0144]$ \\
Random forest (telem) & $0.0758$ $[0.0427,0.1090]$ & $0.0176$ $[0.0099,0.0253]$ \\
\bottomrule
\end{tabular}}
\end{table}


On T1, TabNet attains mean AUPRC $0.0492$ (95\% CI $[0.0213,0.0772]$), below RF ($0.0758$) and far below the manuscript-final ECN-v3 head ($0.1152$). True GraphSAGE attains $0.0152$ (95\% CI $[0.0059,0.0246]$) with $19$ devices and $54$ directed edges in the twin graph---competitive with neither RF nor ECN on this rare-event ranking task. The historical GNN proxy remains at $0.0335$ and must not be interpreted as message-passing performance. On T2, TabNet ($0.0110$) and GraphSAGE ($0.0062$) also trail telem logistic ($0.0380$). These outcomes support reporting deep baselines for completeness while retaining the evidence-selected ECN-v3 / telem-logistic heads.
